% STATUS: rewritten 2026-11-05; every claim is cited.
\section{Lattices and residuated lattices}\label{sec:other-systems:lattices}

Here lattices stand further from a ring than any system above.
A lattice has only an order and two idempotent operations
\cite[p.~309]{XinLiLu2008}.
So a product rule in the old sense has no product to use
\cite[p.~309]{XinLiLu2008}.
For the way in, look to topology:
the
derived set of a subset is a map on a Boolean algebra
\cite[p.~365]{Goldberg1975}.

\begin{definition}[{\cite[p.~365]{Goldberg1975}}]\label{def:other:goldberg-derivation}
Let $R$ be a Boolean algebra. A map $A \mapsto A^{d}$ of $R$ into itself
is a \emph{$d$-operator} if, for all $A, B \in R$,
\textup{(I$_d$)} $(A \cup B)^{d} = A^{d} \cup B^{d}$,
\textup{(II$_d$)} $A^{dd} \subseteq A^{d}$,
\textup{(III$_d$)} $0^{d} = 0$, and
\textup{(IV$_d$)} $1^{d} = 1$.
A map satisfying \textup{(I$_d$)}--\textup{(III$_d$)} is a
\emph{$\delta$-operator}. A Boolean algebra with a $d$-operator (or a
$\delta$-operator) is called a Boolean algebra with a $d$-derivation (or a
$\delta$-derivation).
\end{definition}

The names come from topology. Kuratowski introduced the operator, and if
$R$ is the algebra of all parts of a space without isolated
points, then $A^{d}$ is the derived set of $A$
\cite[p.~365]{Goldberg1975}.

Xin, Li and Lu brought the idea back to lattices in 2008: their
definition reads like the product rule, but it is a rule for the meet
\cite[Definition~3.1, p.~309]{XinLiLu2008}.

\begin{definition}[{\cite[Definition~3.1, p.~309]{XinLiLu2008}}]\label{def:other:lattice-derivation}
Let $L$ be a lattice. A function $d \colon L \to L$ is a \emph{derivation}
on $L$ if
\[
d(x \wedge y) = (dx \wedge y) \vee (x \wedge dy)
\qquad (x, y \in L).
\]
A derivation $d$ on $L$ is \emph{isotone} if $x \le y$ implies $dx \le dy$,
\emph{monomorphic} if it is one-to-one, and \emph{epic} if it is onto
\cite[Definition~3.7, p.~310]{XinLiLu2008}.
\end{definition}

Every derivation on a lattice is decreasing: $dx \le x$ for all $x$
\cite[Proposition~3.6, p.~310]{XinLiLu2008}. The identity is then forced
by any of the usual extra conditions (Theorem~\ref{thm:other:xin-317})
\cite[Theorem~3.17, p.~312]{XinLiLu2008}.\begin{theorem}[{\cite[Theorem~3.17, p.~312]{XinLiLu2008}}]\label{thm:other:xin-317}
Let $d$ be a derivation on a lattice $L$. Then the following are
equivalent: \textup{(1)} $d$ is the identity derivation; \textup{(2)}
$d(x \vee y) = (x \vee dy) \wedge (dx \vee y)$ for all $x, y \in L$;
\textup{(3)} $d$ is a monomorphic derivation; \textup{(4)} $d$ is an epic
derivation.
\end{theorem}

Isotone derivations mark out modular and distributive lattices (Theorem~\ref{thm:other:xin-319})
\cite[Theorem~3.19, p.~313]{XinLiLu2008}.\begin{theorem}[{\cite[Theorem~3.19, p.~313]{XinLiLu2008}}]\label{thm:other:xin-319}
Let $L$ be a modular lattice and $d$ a derivation on $L$. Then the
following are equivalent: \textup{(1)} $d$ is isotone; \textup{(2)}
$d(x \wedge y) = dx \wedge dy$ for all $x, y \in L$; \textup{(3)}
$dx = x$ implies $d(x \vee y) = dx \vee dy$.
\end{theorem}

The next theorem gives the corresponding criterion for a distributive lattice (Theorem~\ref{thm:other:xin-321}).

\begin{theorem}[{\cite[Theorem~3.21, p.~313]{XinLiLu2008}}]\label{thm:other:xin-321}
Let $L$ be a distributive lattice and $d$ a derivation on $L$. Then the
following are equivalent: \textup{(1)} $d$ is isotone; \textup{(2)}
$d(x \wedge y) = dx \wedge dy$ for all $x, y \in L$; \textup{(3)}
$d(x \vee y) = dx \vee dy$ for all $x, y \in L$.
\end{theorem}

{\c{C}}even and {\"O}zt{\"u}rk replaced the second argument by a second
map and opened a wider family of maps (Section~\ref{sec:other-systems:subtraction-algebras} of this
volume returns to them).

\begin{definition}[{\cite[Definition~8, p.~702]{CevenOzturk2008}}]\label{def:other:f-derivation}
Let $L$ be a lattice. A function $d \colon L \to L$ is an
\emph{$f$-derivation} on $L$ if there is a function $f \colon L \to L$
such that
\[
d(x \wedge y) = \bigl(d(x) \wedge f(y)\bigr) \vee \bigl(f(x) \wedge d(y)\bigr)
\qquad (x, y \in L).
\]
If $f$ is the identity function, then $d$ is a derivation on $L$
\cite[p.~702]{CevenOzturk2008}.
\end{definition}

The next proposition records two bounds that every $f$-derivation obeys (Proposition~\ref{prop:other:ceven-prop1}).

\begin{proposition}[{\cite[Proposition~1, pp.~703--704]{CevenOzturk2008}}]\label{prop:other:ceven-prop1}
Let $L$ be a lattice and $d$ an $f$-derivation on $L$. Then, for all
$x, y \in L$, \textup{(a)} $dx \le fx$, and \textup{(b)}
$dx \wedge dy \le d(x \wedge y) \le dx \vee dy$.
\end{proposition}

A residuated lattice carries a product and its residuum, so it has room
for a multiplicative rule.

\begin{definition}[{\cite[Definition~2.1, p.~99]{HeXinZhan2016}}]\label{def:other:residuated-lattice}
A \emph{residuated lattice} is an algebra
$(L, \wedge, \vee, \cdot, \to, 0, 1)$ of type $(2, 2, 2, 2, 0, 0)$
equipped with an order $\le$ such that \textup{(1)}
$(L, \wedge, \vee, 0, 1)$ is a bounded lattice relative to $\le$,
\textup{(2)} $(L, \cdot, 1)$ is a commutative monoid, and \textup{(3)}
$x \cdot y \le z$ if and only if $x \le y \to z$ for all $x, y, z \in L$.
\end{definition}

The next definition states the multiplicative derivations on a residuated lattice.

\begin{definition}[{\cite[Definition~3.1, p.~100]{HeXinZhan2016}}]\label{def:other:multiplicative-derivation}
Let $(L, \wedge, \vee, \cdot, \to, 0, 1)$ be a residuated lattice. A map
$d \colon L \to L$ is a \emph{multiplicative derivation} on $L$ if
\[
d(x \cdot y) = (d(x) \cdot y) \vee (x \cdot d(y))
\qquad (x, y \in L).
\]
\end{definition}

The next definition introduces the ideal derivations.

\begin{definition}[{\cite[Definition~3.6, p.~101]{HeXinZhan2016}}]\label{def:other:ideal-derivation}
A derivation $d$ on a residuated lattice $L$ is \emph{isotone} if
$x \le y$ implies $d(x) \le d(y)$, \emph{contractive} if $d(x) \le x$ for
all $x \in L$, and an \emph{ideal derivation} if it is both.
\end{definition}

The next theorem describes the ideal derivations that are idempotent (Theorem~\ref{thm:other:he-xin-zhan}).

\begin{theorem}[{\cite[Theorem~3.10, p.~102]{HeXinZhan2016}}]\label{thm:other:he-xin-zhan}
Let $(L, \wedge, \vee, \cdot, \to, 0, 1)$ be a residuated lattice and $d$
a derivation on $L$. Then the following are equivalent: \textup{(1)} $d$
is an ideal derivation on $L$ and $d^{2} = d$; \textup{(2)} $d$ satisfies
$d(x) \to d(y) = d(x) \to y$ for all $x, y \in L$.
\end{theorem}

The next section moves to the algebras of logic, where the same rule
appears on a BCI, a BCK or an MV-algebra
\cite[p.~167]{JunXin2004}.

Djouadi and Prade extend formal concept analysis to a possibility-theoretic setting in which four derivation powerset operators appear instead of one, and give minimal requirements on a fuzzy algebra for the closure and opening properties of the antitone and isotone $L$-Galois connections~\cite{WlPossibilityTheoreticExtensionO2011}.

Yang constructs a new class of subspace lattices on an infinite-dimensional Hilbert space, proves that the bounded cohomology groups $H^n(\mathrm{Alg}\,\mathcal{L},B(K))$ of the corresponding lattice algebras with coefficients in $B(K)$ vanish for all $n\geq 1$, and that under stated conditions every derivation of $\mathrm{Alg}\,\mathcal{L}$ into itself is inner~\cite{WlCohomologyGroupsAndLieDerivati2025}.

